\Needspace{12\baselineskip}
\section{References}

\begingroup
\setlength{\parskip}{0.15em}
\titlespacing*{\subsection}{0pt}{1.35ex}{0.45ex}
\begin{multicols}{2}
\raggedcolumns

\subsection*{Jewish origin and reception witnesses}

{\footnotesize
\begin{itemize}[itemsep=0.1em,topsep=0.15em]
\item \emph{Letter of Aristeas}, especially \S\S28--32, 44--50, and 301--318, trans. Herbert T. Andrews, in R.~H. Charles, ed., \emph{The Apocrypha and Pseudepigrapha of the Old Testament}, vol.~2 (Oxford: Clarendon, 1913), \sourceurl{https://ccel.org/c/charles/otpseudepig/aristeas.htm}{public-domain electronic witness}. The work is read as a retrospective Alexandrian Jewish narrative, not as a court archive.
\item Greek prologue to Sirach, especially the translator's arrival under Euergetes, description of prior Greek Scriptures, and statement about translation difficulty, in Benjamin G. Wright's NETS translation, \sourceurl{https://septuaginta.uni-goettingen.de/static/ioscs/pdf/nets/30-sirach-nets.pdf}{official IOSCS PDF}. The wording is paraphrased, not reproduced.
\item Philo, \emph{Life of Moses} II.25--44, trans. C.~D. Yonge, \sourceurl{https://www.earlyjewishwritings.com/text/philo/book25.html}{public-domain English witness}; Josephus, \emph{Antiquities} XII.11--118, \sourceurl{https://www.attalus.org/translate/aj_12a.html}{online English witness}. These establish first-century receptions of the origin tradition, not independent contemporary accounts.
\item Mishnah Megillah 1:8; Babylonian Talmud Megillah 9a; Tractate Soferim 1:7--8, consulted through \sourceurl{https://www.sefaria.org/Mishnah_Megillah.1.8?lang=bi}{Sefaria's mishnaic text}, \sourceurl{https://www.sefaria.org/Megillah.9a?lang=bi}{talmudic text}, and \sourceurl{https://www.sefaria.org/Tractate_Soferim.1.7?lang=bi}{Soferim text}. Their distinct dates and judgments are not harmonized.
\end{itemize}}

\subsection*{Early Christian reception}

{\footnotesize
\begin{itemize}[itemsep=0.1em,topsep=0.15em]
\item Justin Martyr, \emph{Dialogue with Trypho} 71, \sourceurl{https://www.newadvent.org/fathers/01286.htm}{New Advent witness}; Irenaeus, \emph{Against Heresies} III.21.2, \sourceurl{https://www.newadvent.org/fathers/0103321.htm}{New Advent witness}. Their polemical and theological uses do not supply origin-era corroboration.
\item Eusebius, \emph{Ecclesiastical History} VI.16, \sourceurl{https://www.newadvent.org/fathers/250106.htm}{New Advent witness}, for Origen's comparative work and its named versions.
\item Augustine, \emph{City of God} XVIII.42--43, \sourceurl{https://www.newadvent.org/fathers/120118.htm}{New Advent witness}, for mature Latin Christian reception of the translators and meaningful Greek--Hebrew differences.
\end{itemize}}

\subsection*{Manuscripts and material witnesses}

{\footnotesize
\begin{itemize}[itemsep=0.1em,topsep=0.15em]
\item John Rylands Library, Greek Papyrus 458, \sourceurl{https://www.digitalcollections.manchester.ac.uk/view/MS-GREEK-P-00458}{University of Manchester catalog and images}; used for artifact identity, contents, and broad date, not as a complete Torah.
\item Timothy Lee, “The Newly Discovered Fragments of the Greek Minor Prophets Scroll from Na\d{h}al \d{H}ever (8\d{H}evXIIgr) and the Problem of Translation Standardisation,” \emph{Journal of Septuagint and Cognate Studies} 55 (2022), DOI \sourceurl{https://doi.org/10.2143/JSCS.55.0.3291482}{10.2143/JSCS.55.0.3291482}; \sourceurl{https://www.repository.cam.ac.uk/handle/1810/349882}{accepted-manuscript record}.
\item \sourceurl{https://digi.vatlib.it/view/MSS_Vat.gr.1209}{Codex Vaticanus, Vat. gr. 1209, Vatican Library}; \sourceurl{https://www.codexsinaiticus.org/en/codex/content.aspx}{Codex Sinaiticus Project}; and \sourceurl{https://searcharchives.bl.uk/catalog/040-002353500}{Codex Alexandrinus, Royal MSS 1 D V--VIII, British Library}. Official records govern identifiers, custody, surviving contents, and material limits.
\end{itemize}}

\subsection*{Modern historical and textual reconstruction}

{\footnotesize
\begin{itemize}[itemsep=0.1em,topsep=0.15em]
\item Alison G. Salvesen and Timothy Michael Law, eds., \emph{The Oxford Handbook of the Septuagint} (Oxford University Press, 2021), DOI \sourceurl{https://doi.org/10.1093/oxfordhb/9780199665716.001.0001}{10.1093/oxfordhb/9780199665716.001.0001}. Principal chapters consulted were Cameron Boyd-Taylor, “What Is the Septuagint?”; Dirk B\"uchner, “The Pentateuch”; Anneli Aejmelaeus, “The Books of Samuel”; Matthieu Richelle, “Jeremiah and Baruch”; Olivier Munnich, “Daniel, Susanna, Bel and the Dragon”; Maria Gorea, “The Book of Job”; Alison G. Salvesen, “Deuterocanonical and Apocryphal Books”; Eugene Ulrich, “The Scrolls from the Judean Desert”; Siegfried Kreuzer, “Kaige and ‘Theodotion’”; David Lincicum, “Citations in the New Testament”; and Peter J. Gentry, “Origen's Hexapla.”
\item Sylvie Honigman, “Aristeas, Letter of,” \emph{Oxford Research Encyclopedia of Religion} (2016; revised online record), DOI \sourceurl{https://doi.org/10.1093/acrefore/9780199381135.013.741}{10.1093/acrefore/9780199381135.013.741}, for its literary identity, Alexandrian setting, and historiography.
\item Timothy Lee's study of 8\d{H}evXIIgr, above, read alongside Kreuzer's chapter for the distinction between the material scroll, standardizing revision, the kaige category, and later “Theodotion.”
\end{itemize}}

\subsection*{Editions and modern translation}

{\footnotesize
\begin{itemize}[itemsep=0.1em,topsep=0.15em]
\item International Organization for Septuagint and Cognate Studies, \sourceurl{https://septuaginta.uni-goettingen.de/ioscs/editions/}{critical-editions guide}, on Cambridge, Rahlfs--Hanhart, and G\"ottingen editorial aims; \sourceurl{https://adw-goe.de/forschung/abgeschlossene-forschungsprojekte/akademienprogramm/septuaginta-unternehmen/}{G\"ottingen Academy project history}, for the 1 April 1908 start, distinct from the first \emph{editio maior} volume in 1931.
\item Alfred Rahlfs, ed., \emph{Septuaginta}, rev. Robert Hanhart, \emph{editio altera} (German Bible Society, 2006); individual G\"ottingen volumes are preferred where available for book-specific critical work.
\item Albert Pietersma and Benjamin G. Wright, eds., \emph{A New English Translation of the Septuagint and the Other Greek Translations Traditionally Included under That Title} (Oxford University Press, 2007; corrected printing, 2009), \sourceurl{https://septuaginta.uni-goettingen.de/ioscs/nets/edition/}{official electronic edition}. The front matter and Sirach prologue were checked; no NETS biblical text is reproduced here.
\end{itemize}}

\end{multicols}
\endgroup
